% main.tex -- головний файл вашої роботи.
%
% Що тут можна/варто редагувати:
%  - опцію класу нижче: [master] чи [bachelor];
%  - список файлів у chapters/ і appendices/, якщо додаєте/прибираєте
%    розділи чи додатки (просто дописуйте чи коментуйте рядки \input).
%
% Що редагувати НЕ тут: усі особисті дані -- у config.tex.
\documentclass[master]{kpithesis}

% config.tex -- усі дані про вас і вашу роботу в одному місці.
%
% Це ЄДИНИЙ файл, який більшість студентів мають редагувати напряму.
% Він не залежить від факультету чи кафедри -- це просто текстові
% поля, тому шаблон однаково підходить для будь-якої кафедри КПІ:
% просто впишіть свої дані замість прикладів нижче.
%
% Ступінь (бакалавр/магістр) тут НЕ задається -- він вказаний у
% main.tex, в аргументі \documentclass[master]{kpithesis} або
% [bachelor]. Це свідомий поділ: тут -- дані про вас, там -- вибір
% "форми" документа.

% --- Заклад ---
\newcommand{\facultyName}{Повна назва факультету/навчально-наукового інституту}
\newcommand{\departmentName}{Повна назва кафедри}
\newcommand{\departmentHeadName}{Ім'я ПРІЗВИЩЕ}
% Завідувач(-ка) кафедри -- саме той підпис, що стоїть у грифі
% "ДО ЗАХИСТУ ДОПУЩЕНО" на титульній сторінці.

% --- Спеціальність і програма ---
\newcommand{\specialtyCode}{000}
\newcommand{\specialtyName}{Назва спеціальності}
\newcommand{\programmeName}{Назва освітньої програми}
% Для магістерської дисертації важливо, за якою саме програмою
% ви навчаєтесь -- освітньо-професійною чи освітньо-науковою (це
% впливає на формулювання на титульній сторінці). Впишіть одне з двох:
% "освітньо-професійною" або "освітньо-науковою". Для бакалаврської
% роботи це поле не використовується (там завжди ОПП) -- можете
% залишити як є.
\newcommand{\programmeType}{освітньо-професійною}

% --- Тема роботи ---
\newcommand{\thesisTitle}{Тема вашої роботи}
% Для бакалаврської роботи вкажіть одне з двох: "Дипломний проєкт" або
% "Дипломна робота" (перше -- якщо є графічна частина/креслення,
% друге -- якщо ні; для більшості неінженерних спеціальностей це
% "Дипломна робота"). Для магістра це поле не використовується --
% магістерська робота завжди називається "Магістерська дисертація".
\newcommand{\workType}{Дипломна робота}

% --- Автор(-ка) ---
\newcommand{\authorFullName}{Прізвище Ім'я По-батькові}
\newcommand{\authorCourse}{4}
\newcommand{\authorGroup}{XX-00}

% --- Науковий керівник ---
\newcommand{\supervisorFullName}{Прізвище Ім'я По-батькові}
\newcommand{\supervisorRegalia}{посада, науковий ступінь, вчене звання}
% Наприклад: "доцент кафедри ОТ, канд. техн. наук, доцент"

% --- Консультант (за наявності) ---
% Якщо консультанта немає -- просто залиште \consultantFullName
% порожнім (як зараз): титульна сторінка й бланк завдання самі
% пропустять цей рядок.
\newcommand{\consultantFullName}{}
\newcommand{\consultantRegalia}{}
\newcommand{\consultantChapter}{}
% \consultantChapter -- назва розділу, з якого консультували,
% наприклад "економічної частини" -- потрібна лише якщо консультант є.

% --- Рецензент ---
\newcommand{\reviewerFullName}{Прізвище Ім'я По-батькові}
\newcommand{\reviewerRegalia}{посада, науковий ступінь, вчене звання}

% --- Формальності ---
\newcommand{\udcCode}{000.000}
% УДК скаже Вам науковий керівник -- код потрібен
% лише для магістерської дисертації (для бакалаврської не друкується).
\newcommand{\thesisYear}{2026}


\begin{document}

% --- Вступна частина (ДСТУ 3008:2015, розділ 4) ---
\kpithesisTitlePage
\kpithesisTaskPage
% chapters/00-abstract.tex -- Реферат (українською та іноземною мовою).
%
% Вимоги: ДСТУ 3008:2015, п. 4.3; для магістерської дисертації --
% додатково розділ "Реферат" Рекомендацій КПІ (до 500 слів, певна
% послідовність пунктів). Замініть приклад нижче своїм текстом,
% зберігаючи структуру.
\kpithesisUnnumberedChapter{Реферат}

Пояснювальна записка до \workType: \pageref{LastPage}~с.,
\arabic{figure}~рис., \arabic{table}~табл., 2~додатки, 00~джерел.
% ^^^ \pageref{LastPage} працює, лише якщо в кінці документа є мітка
% \label{LastPage} (вона вже додана в main.tex). Кількість джерел
% після підключення бібліографії (Фаза 5) можна порахувати вручну.

\textbf{Об'єкт дослідження} -- \dots

\textbf{Мета роботи} -- \dots

\textbf{Методи дослідження} -- \dots

\textbf{Наукова новизна} (для магістерської дисертації) -- \dots

\textbf{Практичне значення} -- \dots

\bigskip
\noindent\textbf{Ключові слова:} ключове слово 1, ключове слово 2,
ключове слово 3.

\bigskip
\begin{english}
  \textbf{Abstract}

  Explanatory note to the \dots: \pageref{LastPage}~p., \arabic{figure}~fig.,
  \arabic{table}~tables, 2~appendices, 00~references.

  \textbf{Object of research} -- \dots

  \textbf{Purpose} -- \dots

  \textbf{Methods} -- \dots

  \textbf{Scientific novelty} -- \dots

  \textbf{Practical value} -- \dots

  \bigskip
  \noindent\textbf{Keywords:} keyword one, keyword two, keyword three.
\end{english}

\tableofcontents
% chapters/00-abbreviations.tex -- Перелік умовних позначень,
% скорочень і термінів (ДСТУ 3008:2015, п. 4.5 та 7.14).
%
% Потрібен, лише якщо у роботі є нестандартні скорочення/позначення,
% що повторюються 3+ рази (інакше розшифровуйте їх просто в тексті при
% першій згадці). Якщо розділ не потрібен -- не підключайте цей файл
% у main.tex і приберіть \kpithesisChapterConclusion... ні, просто не
% підключайте \input тут.
%
% За абеткою, ліворуч -- скорочення, праворуч -- розшифрування.
\kpithesisUnnumberedChapter{Перелік умовних позначень, скорочень і термінів}

\begin{tabularx}{\textwidth}{@{}l X@{}}
  ПЗ & програмне забезпечення \\
  API & Application Programming Interface, програмний інтерфейс застосунку \\
\end{tabularx}

% ^^^ якщо у роботі немає нестандартних скорочень -- закоментуйте
% рядок вище.

% --- Основна частина (ДСТУ 3008:2015, розділ 5) ---
% chapters/01-intro.tex -- Вступ (ДСТУ 3008:2015, п. 5.1; Рекомендації
% КПІ уточнюють для магістерської дисертації, що обсяг вступу не має
% перевищувати 5% основної частини).
\kpithesisUnnumberedChapter{Вступ}

\textbf{Актуальність теми.} \dots

\textbf{Мета і завдання дослідження.} \dots
Для досягнення мети потрібно розв'язати такі завдання:
\begin{itemize}
  \item \dots;
  \item \dots;
  \item \dots.
\end{itemize}

\textbf{Об'єкт дослідження.} \dots

\textbf{Предмет дослідження.} \dots

\textbf{Методи дослідження.} \dots

\textbf{Наукова новизна одержаних результатів} (обов'язково для
освітньо-наукової програми). \dots

\textbf{Практичне значення одержаних результатів.} \dots

\textbf{Апробація результатів.} \dots

\textbf{Публікації.} \dots

% chapters/02-chapter-01.tex -- перший розділ основної частини.
% Типово тут огляд літератури/сучасного стану проблеми (не більше 20%
% обсягу основної частини -- Рекомендації КПІ).
\chapter{Назва першого розділу}

Текст розділу з посиланням на джерело \cite{makarova2002} та ще одне
\cite{rossokha2015, ifla2012}.

\section{Назва підрозділу}

Текст підрозділу.

\subsection{Назва пункту}

Текст пункту.

\kpithesisChapterConclusion{1}

% chapters/03-chapter-02.tex -- другий розділ основної частини.
\chapter{Назва другого розділу}

Текст розділу.

\section{Назва підрозділу}

Текст підрозділу.

\kpithesisChapterConclusion{2}

% chapters/04-chapter-03.tex -- третій розділ основної частини.
\chapter{Назва третього розділу}

Текст розділу.

\section{Назва підрозділу}

Текст підрозділу.

\kpithesisChapterConclusion{3}

% chapters/05-conclusions.tex -- Висновки (ДСТУ 3008:2015, п. 5.3).
% Тут підсумовують результати УСІЄЇ роботи -- не повторюють текст
% розділів, а формулюють, що з нього випливає.
\kpithesisUnnumberedChapter{Висновки}

У кваліфікаційній роботі отримано такі результати:

\begin{enumerate}
  \item \dots
  \item \dots
  \item \dots
\end{enumerate}


% --- Перелік посилань (ДСТУ 3008:2015, п. 5.5; ДСТУ 8302:2015) ---
\kpithesisUnnumberedChapter{Перелік посилань}
\printbibliography[heading=none]

% --- Додатки ---
% Порядок підключення визначає літеру додатка (перший файл -> "А",
% другий -> "Б" і т.д.) -- не назва файлу.
% appendices/a-appendix.tex -- приклад додатка.
% \kpithesisAppendix{...} сам відкриває нову сторінку, друкує заголовок
% "ДОДАТОК А" (літеру визначає порядок підключення файлів у main.tex,
% не назва файлу!) і скидає лічильники рисунків/таблиць/формул.
\kpithesisAppendix{Лістинг програмного коду}

Текст додатка.

% appendices/b-appendix.tex -- ще один приклад додатка (стане "ДОДАТОК
% Б", бо підключається другим -- див. main.tex).
\kpithesisAppendix{Акт впровадження}

Текст додатка.


\label{LastPage}
\end{document}
